%!TEX root = ./../../thesis-trier.tex

\IN{monoid language}
A submonoid $L \subseteq \Sigma^*$
is also called a \emph{monoid (language)} for short.
The language $L \subseteq \Sigma^*$
is a monoid language if and only if % $LL \subseteq L$, % ne, neutrales element noch drin
$L^* \subseteq L$.
For any $L \subseteq \Sigma^*$, we have $\langle L \rangle_M = L^*$.

For a monoid language $L \subseteq \Sigma^*$
there exists a unique minimal generating set $X$, i.e., a set $X$ minimal for set-inclusion
with $L = X^*$. It is given as 
the set of non-empty words in $L$ that cannot be written as a product of two non-empty words from $L$, namely (see~\cite[p.~4]{lothaire_1997}):
\begin{equation}
	\label{constr_sync:eqn:generating_set_submonoid}
	X = ( L \setminus \{\varepsilon\} ) \setminus ( L \setminus \{\varepsilon\})^2.
\end{equation}
For every set $Y \subseteq \Sigma^*$
such that $Y^* = L$ we have $X \subseteq Y$ (more precisely,
we have $X = ( Y \setminus \{\varepsilon\} ) \setminus ( Y \setminus \{\varepsilon\})^2$).
This implies that~$L$ is generated by a prefix code 
if and only if~$X$ is a prefix code. The next examples
give languages where $X$ is not a code, and where is it a code, but not a prefix code.

\begin{example}
	For $L = \{a,ab,ba\}^*$ we have $X = \{a,ab,ba\}$, which is not a code since $a(ba) = (ab)a$.
\end{example}

\begin{example}
	For $L = \{a,ab\}^*$ we have $X = \{a,ab\}$. 	
\end{example}

\Cref{constr_sync:eqn:generating_set_submonoid} implies that $L$ is regular if and only if $X$ is regular. 

We always have $L \setminus ( L \Sigma^+ \cup \{\varepsilon \} ) \subseteq X$
since $(L \setminus \{\varepsilon\})^2 \subseteq L\Sigma^+$.
In case $X$ is a prefix code, 
the following equalities hold:\footnote{Since codes are defined here as non-empty subsets of $\Sigma^+$,
	for a prefix code~$X$ we have $X^* \ne \{\varepsilon\}$.}
\begin{equation}\label{constr_sync:eqn:prefix_code_gen}
	X = L \setminus ( L \Sigma^+ \cup \{\varepsilon \} )
	= \{ w \in L \setminus \{\varepsilon\} \mid \pref{w} \cap L = \{ w \} \}.
\end{equation} 

Regular languages that are monoids generated by prefix codes
admit an automata-theoretical characterization 
%(that we need in Section~\ref{constr_sync:sec:PSPACE-hardness})
as languages
accepted by PDFAs
whose start state is the only final state (a related statement
is given in \cite[Prop.~3.2]{Eilenberg1974}).

\begin{theorem}[{\cite[Prop.~3.2.5]{DBLP:books/daglib/0025093},\cite[Thm.~3.1]{DBLP:conf/litp/Restivo88}}]
	\label{constr_sync:thm:regular_prefix_codes_and_automata}
	Let $L \subseteq \Sigma^*$.
	Then we have $L = L(\mathcal B)$
	for some PDFA $\automat{B}{P}{\Sigma}{\mu}{p_0}{\{p_0\}}$
	if and only if  $L = \{\varepsilon\}$ or $L$ is a regular language that is a monoid generated by a prefix code that is regular.
\end{theorem}
\begin{proof}	
	We can assume $L \ne \{\varepsilon\}$. Let $L = L(\mathcal B)$
	for a PDFA $\automat{B}{P}{\Sigma}{\mu}{p_0}{\{p_0\}}$.
	Set $X = \{ u \in  L \mid \forall v \in \pref{u} \setminus \{u\} \colon \mu(p_0, v) \ne p_0 \}$.
	Then~$X$ is a regular prefix code (we obtain a PDFA accepting it by removing all outgoing transitions
	from final states in~$\mathcal B$) such that $L = X^*$.
	
	Conversely, let $L = X^*$ where $X$ is a regular prefix code.
	By \cref{constr_sync:prop:aut_prefix_free} 
	we obtain a trim PDFA $\automat{B}{P}{\Sigma}{\mu}{p_0}{\{q\}}$
	such that $q$ has no outgoing transitions.
	Thus we can merge the states $q_0$ and~$q$ to a single
	state $p$ to obtain a new PDFA where all transitions going out of~$p_0$ go out of this new state
	and all transitions ending in some state from $\{q, p_0\}$
	now end at~$p$.
	This new PDFA with start state~$p$ and final state set~$\{p\}$
	accepts~$L$.
\end{proof}

\Cref{constr_sync:thm:regular_prefix_codes_and_automata} has been extended in~\cite[Thm.~3.2]{DBLP:conf/litp/Restivo88}
to characterize submonoids that are regular languages as precisely those languages $L \subseteq \Sigma^*$
such that $L = \{ u \in \Sigma^* \mid \applytrans{\mu}{S}{u} = S \}$ for a PDFA $\automat{B}{P}{\Sigma}{\mu}{p_0}{F}$
and $S \subseteq P$.

\IN{literal automaton}
\begin{remark} \label{syncmax:rem:general-aut-prefix-code}
	The construction from the proof of
	\crefitem{constr_sync:prop:aut_prefix_free}{constr_sync:prop:aut_prefix_free:item:min}
	(used in the proof of \cref{constr_sync:thm:regular_prefix_codes_and_automata})
	is closely related to the rooted tree from
	\cref{syncmax:rem:root-tree-prefix-code}.
	In fact, this rooted tree can be considered as a (possibly infinite)
	PDFA accepting the prefix code.
	This PDFA is called the \emph{literal automaton}
	of the prefix code in~\cite[Sect.~3.2]{DBLP:books/daglib/0025093}.
	See also~\cite[Sect.~1.3 \& 4.3.8]{DBLP:phd/hal/Ryzhikov20},
	where this construction is used for finite maximal prefix codes.
	In the construction from
	\crefitem{constr_sync:prop:aut_prefix_free}{constr_sync:prop:aut_prefix_free:item:min},
	we additionally incorporate the Myhill--Nerode relation to obtain a finite automaton whenever the prefix code is regular.
\end{remark}

Next, we show that monoid languages in $u^*$ for some $u \in \Sigma^*$ are always regular.

\IN{Frobenius coin problem}
\begin{proposition}\label{constrsync:thm:sparsemonoid-regular}
	Le $L \subseteq u^*$ be a monoid language for some $u \in \Sigma^*$.
	Then there exist $d \ge 1$ and a multiple of $n$ of~$d$
	such that $L = u^n(u^d)^* \cup F$
	for a finite set $F \subseteq \{ u \}^{< n} \cap (u^d)^*$. 
	In particular, $L$ is regular.
\end{proposition}
\begin{proof}
	We can assume $L$ is infinite.
	Let $d = \min\{ \gcd(i_1, \ldots, i_r) \mid \{ u^{i_1}, \ldots, u^{i_r} \} \subseteq L\}$.		
	Then there exist numbers $i_1, \ldots, i_r \ge 0$ with $\{ u^{i_1}, \ldots, u^{i_r} \} \subseteq L$
	such that $d = \gcd( i_1, \ldots, i_r )$.
	Let $i \ge 0$ with $u^i \in L$.
	As $d = \gcd( a_1, \ldots, a_r, i )$ 
	the number~$d$ divides~$i$.
	This yields $L \subseteq (u^d)^*$.	
	By B\'{e}zout's identity 
	there exists a number $n \ge 0$ divisible by $d$ such that every
	number $m$ divisible by $d$
	with $m \ge n$ can be written as $m = a_1 i_1 + \cdots + a_r i_r$
	for numbers $a_1, \ldots, a_r \ge 0$.\footnote{We note that the \emph{Frobenius coin problem}
		asks for the smallest $n$ for given $i_1, \ldots, i_r$ when~$d = 1$, see~\cite{Alfonsis05}.}
	So for $u^m \in u^n(u^d)^*$ we have
	$u^m \in  (u^{i_1})^* \cdots (u^{i_r})^* \subseteq L$
	and we get $u^n (u^d)^* \subseteq L$.
	As $L \subseteq (u^d)^*$ and $n$ is a multiple of $d$, the language
	$F = L \setminus u^n (u^d)^*$ is finite and $F \subseteq \{u\}^{< n}$.
	Hence, $L = a^n (u^d)^* \cup F$.
\end{proof}

The next corollary generalizes the well-known folklore fact
that the Kleene star of any unary language is always regular.
Via the mapping $n \mapsto u^n$ for a fixed $u \in \Sigma^*$ 
it is in fact a reformulation of the statement
that a submonoid of~$\mathbb N_0$
is ultimately periodic~\cite[Prop.~2.2]{RosalesGarciaSanchez2009}.

\begin{corollary}
	Let $L \subseteq u^*$. Then $L^*$ is a regular language.
\end{corollary}
\begin{proof}
	As $L^*$ is a monoid and since $L^* \subseteq (u^*)^* = u^*$
	by \cref{prelim:langeqn:Kleenestar-idemp},
	the claim follows from 
	\cref{constrsync:thm:monoid-sparseregular-or-exp}.  
\end{proof} 

%https://www.google.com/search?q=kleene+star+of+unary+language+always+regular&rlz=1C1GCEU_deDE823DE823&oq=kleene+star+of+unary+language+always+regular&aqs=chrome..69i57j33i160.12614j0j7&sourceid=chrome&ie=UTF-8
%https://cs.stackexchange.com/questions/21765/kleene-star-of-an-infinite-unary-language-always-yields-a-regular-language
%https://cs.stackexchange.com/questions/10013/if-l-is-a-subset-of-0-then-how-can-we-show-that-l-is-regular
%https://cs.stackexchange.com/questions/164/what-are-the-possible-sets-of-word-lengths-in-a-regular-language/175#175
%https://mathworld.wolfram.com/FrobeniusNumber.html
%https://www.youtube.com/watch?v=H-GYBjDpT6U

\IN{semigroup language}
\begin{remark}
	A \emph{semigroup language} is a language $L \subseteq \Sigma^*$
	such that $L^+ \subseteq L$. Equivalently, $L$ forms a subsemigroup of~$\Sigma^*$.
	Every monoid language
	is also a semigroup language, but not vice versa. For example, $aaaa^*$
	is a semigroup language that is not a monoid language.
	However, a language $L \subseteq \Sigma^*$ is a semigroup language
	if and only if $L \cup \{\varepsilon\}$ is a monoid language.
	Similar statements can be given for semigroup languages and 
	our later results on the constrained synchronization problem w.r.t.~monoid
	languages are easily adaptable and also hold for semigroup languages.
\end{remark}